\begin{abstract}
Event clustering and event-dependent inhibition can coexist in a point process, yet a plausible physical construction does not ensure that their contributions can be distinguished from an observed catalogue. This manuscript studies that distinction for marked earthquake models with relaxation memory. It separates an autonomous physical loading model from a finite-horizon diagnostic with an independently specified background, and distinguishes finite positive exponential mixtures from exact fractional memory.

The mathematical analysis gives restricted first-moment stability and stationary-balance results for finite-mode processes, explains how these can coexist with infinite stationary second moments, and identifies, in no-feedback submodels, explicit obstructions to an unrestricted exact-memory extension. Structural identification at positive inhibition is established only under the stated observation and mark assumptions; at zero inhibition a separately varied stress-memory kernel disappears from the diagnostic observation law. Conditional likelihood identities and positive-inhibition response calculations provide useful finite-history tools without supplying general boundary asymptotics or a universal catalogue-size requirement.

Saved controlled experiments distinguish observation-law separability from the performance of procedures that estimate nuisance structure. A known-law test detected all 150 saved alternatives, with 50 rejections among 1,000 separate null queries. Detection was 123/150 when background level was estimated with trend supplied, and 68/150 when level and trend were both estimated. These are exploratory, procedure-specific calibrated comparisons on reused histories, not a causal decomposition or a guarantee of composite-test performance. In a reconstructed 913-event New Zealand catalogue, the declared shared-inhibition recipe did not improve predictive score over its noninhibited Mittag--Leffler comparator in any of three test periods; most of its aggregate disadvantage occurred in the first period. The results support a bounded mathematical and methodological contribution together with negative predictive evidence about the tested recipes. They establish neither the presence nor the absence of a shared physical stress-memory mechanism.
\end{abstract}
